\begin{textsample}{POS Dim 2 – vem\_ed\_21 – Score 9.00 – vem\_ed\_21/t106.txt}  \label{ex:f2_pos_167}
Oficinas Desafios do \textbf{treinamento} e \textbf{sustentabilidade} dos projetos Regiane Souza Camargo Moreira (equipe dos PCIs / Cesu) ofereceu a oficina"The Challenge of VE Training Everyone, Everywhere, All at Once!"(O desafio de treinar Intercâmbios Virtuais para todos, em qualquer lugar, ao mesmo tempo) com Gisselle Morales, Jorge Membrillo, Monserrat Balandrano (Tecnológico de Monterrey / México) e Alonso Cruz (Florida International University / EUA).
Para ter sucesso, os coordenadores de Intercâmbios Virtuais \textbf{devem} garantir que os professores possuam competências de comunicação intercultural, pensamento crítico, cidadania global, colaboração multicultural, adaptação a ambientes de trabalho virtuais e uso de tecnologias.
Após conhecer as práticas do Tecnológico de Monterrey, da Florida International University e do Centro Paula Souza, os ouvintes, reunidos em grupos, debateram desafios relativos a suas experiências de Intercâmbios Virtuais.
Foram convidados a \textbf{responder} a uma pesquisa e os resultados \textbf{devem} ser \textbf{divulgados} futuramente em publicação de acesso aberto.
Osvaldo Succi Junior \textbf{conduziu} a oficina"Scaling COIL on a budget"(Fazendo os projetos COIL crescer com pouco dinheiro) com Eva Haug (Amsterdam University of Applied Sciences / Holanda) e a participação remota de Alicia Betts (Universitat de Girona / Espanha) e Marina Vives (Universitat Rovira i Virgili / Espanha).
Todos são coordenadores de Intercâmbios Virtuais em suas instituições e compartilharam estratégias para a \textbf{sustentabilidade} dos projetos COIL, como por exemplo: capacitação e desenvolvimento profissional, \textbf{reciclagem} de projetos COIL, \textbf{oferta} de COILs padronizados, parcerias sustentáveis, compartilhamento de boas práticas e divulgação por meio de"\textbf{embaixadores} COIL".

% matched lemmas: conduzir, dever, divulgar, embaixadores, oferta, reciclagem, responder, sustentabilidade, treinamento
\end{textsample}
